\subsection{虚数解型隣接3項間漸化式}\label{節：虚数解型隣接3項間漸化式}
  \BookFigComplexRoots

  今までの隣接3\,項間漸化式(\sectionref{節：隣接3項間漸化式型})では,特性方程式の解$\alpha,\beta$が異なれば$a_n=A\alpha^{n-1}+B\beta^{n-1}$という式自体は虚数解の場合でもそのまま成り立つ.
  ここでは,特性方程式が虚数解をもつ場合に,この式を三角関数を用いた実数の形に書き直す方法を考える.

  本項では,実数解で使った「特性根の冪を組み合わせる」流れを複素数へ広げる.
  \sectionref*{節：解を組み合わせる見方}で学んだ線形結合と初期条件の一意性を使い,
  実数の形への書き直しを確かめる.その後に,補助数列と加法定理から同じ形を見る.
  実数の係数$p,q$について,まず,
  \[
    a_{n+2}+pa_{n+1}+qa_n=0
  \]
  を考える.特性方程式は
  \[
    x^2+px+q=0
  \]
  である.

  これまで,特性方程式の解が異なる実数
  \[
    x=\alpha,\beta
  \]
  をもつ場合には,
  \[
    a_n=A\alpha^{n-1}+B\beta^{n-1}
  \]
  という形で考えた.本項の非実の根は零でないので,指数を一つずらした分を係数に吸収できる.
  以下では指数を$n$にした表し方も使う.

  では,特性方程式が実数解をもたず,
  \[
    x=\alpha,\overline{\alpha}
  \]
  という互いに共役な解をもつ場合には,一般項はどのような形になるだろうか.

  ここで「二つの冪で解全体を表せる」という理由を,前の共通説明につなげておこう.
  $p^2-4q<0$だから$q>0$であり,$\alpha,\overline\alpha$は異なる非零の根である.
  $u_n=\alpha^{n-1}$を元の式へ代入すると,
  \[
    u_{n+2}+pu_{n+1}+qu_n
    =\alpha^{n-1}(\alpha^2+p\alpha+q)=0
  \]
  となる.$v_n=\overline\alpha^{\,n-1}$も同じ式を満たす.
  それらの初めの二項は$(1,\alpha)$と$(1,\overline\alpha)$である.
  $Cu_n+Dv_n=0$なら$n=1,2$から
  \[
    C+D=0,\qquad C\alpha+D\overline\alpha=0
  \]
  を得る.$\alpha\ne\overline\alpha$なので$C=D=0$であり,二つの解は一次独立である.
  また,どんな初期条件もこの二つの係数で一意に合わせられる.
  従って,\sectionref*{節：一次独立・基底・次元}で述べた2次元の解全体の基底になっている.
  ここでは係数$C,D$も複素数として考えている.
  初めの二項を合わせれば,同じ更新からすべての項の一致が分かる.
  非実の根でも,実数解のときと同じ理由で冪の線形結合を使えるのである.

  例えば,
  \[
    x^2-2x+2=0
  \]
  を考えると,
  \[
    x=1\pm i
  \]
  である.したがって,漸化式
  \[
    a_{n+2}=2a_{n+1}-2a_n
  \]
  の特性方程式は虚数解をもつ.

  したがって,この例の一般項を
  \[
    a_n=A(1+i)^n+B(1-i)^n
  \]
  と書くことができる.しかし,\,$a_n$が実数列である場合,このままでは一般項の実数性が分かりにくい.そこで,複素数の極形式を利用する.

  一般に,特性方程式の解を
  \[
    \alpha=\rho(\cos\theta+i\sin\theta)=\rho e^{i\theta}
  \]
  とすると,共役解は
  \[
    \overline{\alpha}=\rho(\cos\theta-i\sin\theta)=\rho e^{-i\theta}
  \]
  である.したがって,
  \[
    \alpha^n=\rho^n e^{in\theta}
    =\rho^n(\cos n\theta+i\sin n\theta)
  \]
  および
  \[
    \overline{\alpha}^{\,n}
    =\rho^n e^{-in\theta}
    =\rho^n(\cos n\theta-i\sin n\theta)
  \]
  となる.

  ここで,$a_n$が実数列なら,共役を取っても値が変わらないので,
  \[
    A\alpha^n+B\overline\alpha^{\,n}
    =\overline B\alpha^n+\overline A\,\overline\alpha^{\,n}
  \]
  である.先ほど確かめた一次独立性から,係数の表し方は一意だから,
  $B=\overline A$となる.実数の解を選ぶ条件が,係数を共役にすることとして現れた.
  実際,\,$A=\dfrac{C-iD}{2},\ B=\dfrac{C+iD}{2}$(\,$C,D$は実数)とおくと,
\BookMathLayout{\begin{align*}
    A\alpha^n+B\overline{\alpha}^{\,n}
    &=\frac{C-iD}{2}\rho^n(\cos n\theta+i\sin n\theta)\\
    &\qquad+\frac{C+iD}{2}\rho^n(\cos n\theta-i\sin n\theta)\\
    &=\rho^n(C\cos n\theta+D\sin n\theta)
  \end{align*}}{
\[
\begin{aligned}
A\alpha^n+B\overline\alpha^{\,n}\\
{}=\frac{C-iD}{2}\rho^n(\cos n\theta+i\sin n\theta)\\
\quad+\frac{C+iD}{2}\rho^n(\cos n\theta-i\sin n\theta)\\
{}=\rho^n(C\cos n\theta+D\sin n\theta)
\end{aligned}
\]
}
  のように,虚部が打ち消し合って
  \[
    a_n=\rho^n(C\cos n\theta+D\sin n\theta)
  \]
  という実数の形に書き直せることが分かる.

  したがって,虚数解をもつ隣接3\,項間漸化式では,
  \[
    a_n=\rho^n(C\cos n\theta+D\sin n\theta)
  \]
  という形が基本的な一般項となる.ここで,\,$\rho$は特性根の絶対値,\,$\theta$はその偏角である.

  では,一般項
  \[
    a_n=\rho^n(C\cos n\theta+D\sin n\theta)
  \]
  に含まれる定数$C,D$をどのように決定すればよいだろうか.

  これまでの漸化式と同様に,初期条件を代入することで$C,D$を決定することができる.例えば,\,$a_1,a_2$が与えられているとする.一般項に$n=1,2$を代入すると,
  \[
    a_1=\rho(C\cos\theta+D\sin\theta)
  \]
  および
  \[
    a_2=\rho^2(C\cos2\theta+D\sin2\theta)
  \]
  を得る.

  これは$C,D$についての連立方程式である.したがって,この2つの方程式を解けば$C,D$が決定される.

  特に,$a_0,a_1$から始め,漸化式を$n\ge0$で使う場合には,より簡単になる.一般項に$n=0$を代入すると,
  \[
    a_0=C
  \]
  となる.したがって,
  \[
    C=a_0
  \]
  が直ちに得られる.

  さらに,\,$n=1$を代入すると,
  \[
    a_1=\rho(C\cos\theta+D\sin\theta)
  \]
  であるから,
  \[
    D=
    \cfrac{\frac{a_1}{\rho}-C\cos\theta}{\sin\theta}
  \]
  となる.ここで,\,$\sin\theta\neq0$であることに注意する.特性根が実数ではなく虚部をもつ場合には,\,$\theta$は$0$や$\pi$ではないため,
  \[
    \sin\theta\neq0
  \]
  が成り立つ.

  したがって,\,$a_0,a_1$が与えられている場合には,
  \[
    C=a_0,\qquad
    D=\cfrac{\frac{a_1}{\rho}-a_0\cos\theta}{\sin\theta}
  \]
  として$C,D$を求めることができる.

  例えば,先ほどの
  \[
    a_{n+2}=2a_{n+1}-2a_n
  \]
  を考える.特性根は
  \[
    1+i
    =\sqrt{2}
    \left(
      \cos\frac{\pi}{4}
      +i\sin\frac{\pi}{4}
    \right)
  \]
  であるから,
  \[
    \rho=\sqrt{2},\qquad
    \theta=\frac{\pi}{4}
  \]
  である.したがって一般項は
  \[
    a_n=(\sqrt{2})^n
    \left(
      C\cos\frac{n\pi}{4}
      +D\sin\frac{n\pi}{4}
    \right)
  \]
  となる.

  例えば初期条件
  \[
    a_0=1,\qquad a_1=2
  \]
  が与えられているとする.まず,
  \[
    C=a_0=1
  \]
  である.また,
  \[
    2=\sqrt{2}
    \left(
      C\cos\frac{\pi}{4}
      +D\sin\frac{\pi}{4}
    \right)
  \]
  であるから,
  \[
    2=\sqrt{2}
    \left(
      \frac{1}{\sqrt{2}}
      +\frac{D}{\sqrt{2}}
    \right)
    =1+D
  \]
  より,
  \[
    D=1
  \]
  を得る.

  よって,この場合の一般項は
  \[
    a_n=(\sqrt{2})^n
    \left(
      \cos\frac{n\pi}{4}
      +\sin\frac{n\pi}{4}
    \right)
  \]
  となる.

  このように,虚数解をもつ場合でも,基本的な流れは実数解の場合と変わらない.まず特性根から$\rho,\theta$を求め,
  \[
    a_n=\rho^n(C\cos n\theta+D\sin n\theta)
  \]
  とおき,初期条件を代入して$C,D$を決定すればよい.

  \noindent\textbf{【発想のつながり：倍率をそろえると加法定理が現れる】}\par
  特性根を極形式で書いたのは,一歩ごとの倍率と回転を読み取るためだった.
  解と係数の関係から$p=-2\rho\cos\theta,q=\rho^2$なので,元の式は
  \[
    a_{n+2}-2\rho\cos\theta\,a_{n+1}+\rho^2a_n=0
  \]
  である.今度はこの係数の形から補助数列を設計してみよう.
  初項$a_1$に合わせて指数を$n-1$にし,以下では別の係数$c,d$を使う.
  係数に$\rho$と$\rho^2$があるので,添字一つにつき$\rho$倍される量を基準にする.
  \sectionref*{節：補助数列の設計}の正規化を使い,
  \[
    b_n=\frac{a_n}{\rho^{n-1}}
  \]
  と置く.元の式を$\rho^{n+1}$で割ると
  \[
    b_{n+2}-2\cos\theta\,b_{n+1}+b_n=0
  \]
  となる.大きさの倍率を除いたことで,加法定理と同じ関係が残った.
  実際,
  \[
    \cos(x+\theta)+\cos(x-\theta)=2\cos\theta\cos x
  \]
  と,余弦を正弦に替えた同じ公式に$x=n\theta$を入れると,
  \[
    u_n=\cos((n-1)\theta),\qquad
    v_n=\sin((n-1)\theta)
  \]
  はどちらも$w_{n+2}-2\cos\theta\,w_{n+1}+w_n=0$を満たすと分かる.
  \sectionref*{節：線形結合}の重ね合わせの原理から,$cu_n+dv_n$も同じ関係を満たす.
  初めの二項は$(1,\cos\theta)$と$(0,\sin\theta)$であり,$\sin\theta\ne0$なので一次独立である.
  先ほどの複素数の冪に代わり,正弦・余弦が実数の解全体の基底になる.
  そこで
  \[
    b_n=c\cos((n-1)\theta)+d\sin((n-1)\theta)
  \]
  を候補にしよう.

  \noindent\textbf{【初期条件を合わせ,すべての項で確かめる】}\par
  候補の初めの二項は$c$と$c\cos\theta+d\sin\theta$である.
  元の初期条件から$b_1=a_1,b_2=\frac{a_2}{\rho}$なので,
  \[
    c=a_1,\qquad
    d=\cfrac{\frac{a_2}{\rho}-a_1\cos\theta}{\sin\theta}
  \]
  と選べば一致する.$\sin\theta\ne0$だから,この二つの係数は必ず決められる.
  候補も元の$b_n$も,直前の二項から同じ式で次の項を作る.
  従って初めの二項が一致すれば第三項も一致し,同じことを繰り返すとすべての項で一致する.
  初期条件と加法定理だけで,この候補が一般項であることを確かめられた.
  元へ戻すと,
  \BookMathLayout{\[
    a_n=\rho^{n-1}\{c\cos((n-1)\theta)+d\sin((n-1)\theta)\}
  \]}{\[
    \begin{aligned}
      a_n=\rho^{n-1}\bigl\{&c\cos((n-1)\theta)\\
                           &+d\sin((n-1)\theta)\bigr\}
    \end{aligned}
  \]}
  を得る.倍率をそろえ,加法定理に合う量を作り,初めの二項を合わせるのが解法の流れである.

  \noindent\textbf{【振動と周期性を区別する】}\par
  $C=D=0$なら零数列であり,以下では少なくとも一方が零でないとする.
  三角関数の部分には振動する形が現れるが,整数の添字で見た数列が周期的になるとは限らない.
  $\rho^{n}$は全体の大きさの倍率なので,
  \[
    \begin{cases}
      0<\rho<1 & \text{振幅が減衰する},\\
      \rho=1 & \text{振幅が一定である},\\
      \rho>1 & \text{振幅が増大する}
    \end{cases}
  \]
  と区別できる.
  $\rho=1$で$\frac{\theta}{2\pi}$が有理数なら,正整数$T$について$T\theta$を$2\pi$の整数倍にでき,
  加法定理から$a_{n+T}=a_n$となる.
  非零の解については,逆も成り立つ.
  複素数の冪の式で周期$T$を仮定すると,
  \[
    A(\alpha^T-1)\alpha^{n}
    +\overline A(\overline\alpha^{\,T}-1)\overline\alpha^{\,n}=0
  \]
  を得る.$n=1,2$の二つの式を使うと,$\alpha\ne\overline\alpha$だから両方の係数が零になる.
  非零の実数解では$A,\overline A$も零でないので,$\alpha^T=\overline\alpha^{\,T}=1$である.
  従って周期性には$\rho=1$と$\frac{\theta}{2\pi}\in\mathbb Q$の両方を確認する.
  三角関数があるというだけで周期性を結論しないようにしよう.

  次の数列の一般項を求めてみよう.

  \noindent \bookblocklink{example-complex-1}{problem}{answer}{【例題】}
  \begin{enumerate}[label=(\arabic*),parsep=3pt,format=\bookitemlink{example-complex-1}{problem}{answer}]
    \item $\displaystyle a_1=1,\,a_2=0,\quad a_{n+2}=-a_n\quad(n\ge1)$
\item \BookMathLayout{$\displaystyle a_1=0,\,a_2=1,\quad a_{n+2}=2a_{n+1}-2a_n\quad(n\ge1)$}{$\displaystyle a_1=0,\,a_2=1$\\ $\displaystyle a_{n+2}=2a_{n+1}-2a_n$\\ $\displaystyle(n\ge1)$}
    \item $\displaystyle a_1=1,\,a_2=2,\quad a_{n+2}=-4a_n\quad(n\ge1)$
\item \BookMathLayout{$\displaystyle a_1=1,\,a_2=0,\quad a_{n+2}=2a_{n+1}-4a_n\quad(n\ge1)$}{$\displaystyle a_1=1,\,a_2=0$\\ $\displaystyle a_{n+2}=2a_{n+1}-4a_n$\\ $\displaystyle(n\ge1)$}
  \end{enumerate}

  \noindent\bookblocklink{example-complex-1}{hint}{problem}{【方針】}
  \begin{enumerate}[label=(\arabic*),parsep=3pt,format=\bookitemlink{example-complex-1}{hint}{problem}]
    \item 初項から数項を計算して周期性を見つけ,一般項を予想します.その後,数学的帰納法で証明します.
    \item 特性根を極形式で表し,$\rho=\sqrt2,\,\theta=\dfrac{\pi}{4}$を用いて一般項を置きます.
    \item 特性根は$\pm2i$です.$\rho=2,\,\theta=\dfrac{\pi}{2}$として初期条件を代入します.
    \item 特性根を極形式で表し,$\rho=2,\,\theta=\dfrac{\pi}{3}$として係数を決定します.
  \end{enumerate}

  \noindent\bookblocklink{example-complex-1}{answer}{problem}{【解答】}
  \begin{enumerate}[label=(\arabic*),parsep=3pt,format=\bookitemlink{example-complex-1}{answer}{problem}]
    \item $\displaystyle a_n=\cos\frac{(n-1)\pi}{2}$
    \item $\displaystyle a_n=(\sqrt2)^{n-1}\sin\frac{(n-1)\pi}{4}$
    \item $\displaystyle a_n=2^{n-1}\left(\cos\frac{(n-1)\pi}{2}+\sin\frac{(n-1)\pi}{2}\right)$
\item \BookMathLayout{$\displaystyle a_n=2^{n-1}\left(\cos\frac{(n-1)\pi}{3}-\frac1{\sqrt3}\sin\frac{(n-1)\pi}{3}\right)$}{$\displaystyle a_n=2^{n-1}\Bigl(\cos\frac{(n-1)\pi}{3}$\\
$\displaystyle\hspace{3zw}-\frac1{\sqrt3}\sin\frac{(n-1)\pi}{3}\Bigr)$}
  \end{enumerate}

  \noindent\bookblocklink{example-complex-1}{solution}{problem}{【解法】}
  \begin{enumerate}[label=(\arabic*),parsep=3pt,format=\bookitemlink{example-complex-1}{solution}{problem}]
    \item 初項から数項を計算すると,
          \begin{align*}
            &a_1=1,\quad a_2=0,\quad a_3=-1,\quad a_4=0,\\
            &a_5=1,\quad a_6=0,\quad a_7=-1,\quad a_8=0
          \end{align*}
          となる.この並びから,
          \[
            a_n=\cos\frac{(n-1)\pi}{2}
          \]
          と予想できる.これを数学的帰納法で証明する.
          $n=1,2$では
          \[
            a_1=1=\cos0,\qquad a_2=0=\cos\frac{\pi}{2}
          \]
          となり,主張が成り立つ.
          ある$k\ge1$で,$n=k,k+1$に対して
          \[
            a_k=\cos\frac{(k-1)\pi}{2},\qquad
            a_{k+1}=\cos\frac{k\pi}{2}
          \]
          が成り立つと仮定する.漸化式と$-\cos x=\cos(x+\pi)$より,
          \[
            \begin{aligned}
              a_{k+2}
              &=-a_k\\
              &=-\cos\frac{(k-1)\pi}{2}\\
              &=\cos\left(\frac{(k-1)\pi}{2}+\pi\right)\\
              &=\cos\frac{(k+1)\pi}{2}
            \end{aligned}
          \]
          である.したがって,$n=k+1,k+2$でも主張が成り立つ.初めの2項で成り立つことと合わせ,数学的帰納法により
          \[
            a_n=\cos\frac{(n-1)\pi}{2}
          \]
          がすべての$n\ge1$で成り立つ.
    \item 特性方程式は
          \[
            x^2-2x+2=0
          \]
          であり,特性根は$1\pm i$である.したがって
          \[
            \rho=\sqrt2,\qquad\theta=\frac{\pi}{4}
          \]
          として,
\BookMathLayout{\[
            a_n=(\sqrt2)^{n-1}
            \left(
              C\cos\frac{(n-1)\pi}{4}
              +D\sin\frac{(n-1)\pi}{4}
            \right)
          \]}{
\[
\begin{aligned}
a_n&=(\sqrt2)^{n-1}\Bigl(C\cos\frac{(n-1)\pi}{4}\\
&\hspace{4zw}+D\sin\frac{(n-1)\pi}{4}\Bigr)
\end{aligned}
\]
}
          とおける.$a_1=0$より$C=0$であり,$a_2=1$より$D=1$である.よって,
          \[
            a_n=(\sqrt2)^{n-1}\sin\frac{(n-1)\pi}{4}
          \]
          を得る.
    \item 特性方程式は$x^2+4=0$であり,特性根は$2i,-2i$である.したがって
          \[
            a_n=2^{n-1}
            \left(
              C\cos\frac{(n-1)\pi}{2}
              +D\sin\frac{(n-1)\pi}{2}
            \right)
          \]
          とおける.$a_1=1$より$C=1$であり,$a_2=2$より$2D=2$,すなわち$D=1$である.よって,
          \[
            a_n=2^{n-1}
            \left(
              \cos\frac{(n-1)\pi}{2}
              +\sin\frac{(n-1)\pi}{2}
            \right)
          \]
          を得る.
    \item 特性方程式は
          \[
            x^2-2x+4=0
          \]
          であり,特性根は$1\pm\sqrt3\,i$である.したがって
          \[
            \rho=2,\qquad\theta=\frac{\pi}{3}
          \]
          として,
          \[
            a_n=2^{n-1}
            \left(
              C\cos\frac{(n-1)\pi}{3}
              +D\sin\frac{(n-1)\pi}{3}
            \right)
          \]
          とおける.$a_1=1$より$C=1$であり,$a_2=0$より
          \[
            0=2\left(\frac12+D\frac{\sqrt3}{2}\right)
          \]
          となるので$D=-\dfrac1{\sqrt3}$である.よって,
          \[
            a_n=2^{n-1}
            \left(
              \cos\frac{(n-1)\pi}{3}
              -\frac1{\sqrt3}\sin\frac{(n-1)\pi}{3}
            \right)
          \]
          を得る.
  \end{enumerate}
